\begin{definition}[non-decreasing-$\gamma_k$-Littlestone tree]\label{def:monotone-gamma-littlestone}
Fix a non-decreasing sequence $(\gamma_k)_{k\ge1}$ with $\gamma_k\to\infty$.
A \emph{non-decreasing-$\gamma_k$-Littlestone tree for $\H$}
(a metric-loss variant of the classical Littlestone
tree~\citep{littlestone_learning_1988}; see
also~\citet[Definition~1.7]{bousquet_theory_2021})
is a rooted full
binary tree whose internal nodes at depth $k$ are labeled by instances
$x_{k,i}\in\X$, and whose two outgoing edges from each such node are labeled by
$y_{k,i,1},y_{k,i,2}\in\Y$ satisfying $\ell(y_{k,i,1},y_{k,i,2})\ge\gamma_k$,
with the property that every \emph{finite} root-to-leaf path is realizable by
some $h\in\H$ (i.e., $h$ matches the edge labels along that finite path).
\end{definition}

\begin{definition}[realizable non-decreasing-$\gamma_k$-Littlestone tree]\label{def:realizable-monotone-gamma-littlestone}
An \emph{infinite non-decreasing-$\gamma_k$-Littlestone tree} is
a tree of infinite depth satisfying the conditions of
Definition~\ref{def:monotone-gamma-littlestone} at every finite depth $k$.
Such an infinite tree is \emph{realizable} if,
in addition to finite-prefix realizability, every infinite root-to-leaf path
has its entire labeling realized by a \emph{single} hypothesis in~$\H$ (i.e., for
each infinite path there exists $h\in\H$ matching all edge labels along that
path).
\end{definition}

\begin{lemma}[Bridging lemma: finite-prefix realizability implies
  infinite-path realizability]\label{lem:bridging}
Suppose $\H$ satisfies the measurability condition of
\citet[Definition~3.3]{bousquet_theory_2021} with a \textbf{compact} parameter space
$\Theta$ and a function $h\colon\Theta\times\X\to\Y$ that is \textbf{continuous in~$\theta$} (i.e., $\theta\mapsto h(\theta,x)$ is continuous for each fixed $x\in\X$).
Then every infinite non-decreasing-$\gamma_k$-Littlestone tree
(Definition~\ref{def:monotone-gamma-littlestone}) is automatically realizable
(Definition~\ref{def:realizable-monotone-gamma-littlestone}).
\end{lemma}
\begin{proof}
The proof follows from the finite-intersection property
of compact spaces; see Appendix~\ref{app:compact-parameterization}.
\end{proof}

\begin{remark}\label{rem:compact-param}
Lemma~\ref{lem:bridging} requires the compact-parameterization
assumption (compact~$\Theta$, $h$ continuous in~$\theta$); we discuss in
Appendix~\ref{app:compact-parameterization} why this assumption is
necessary in general, give a concrete example where the two tree
notions differ without it, and argue that the assumption is mild and
covers virtually all settings of practical interest.
\end{remark}

\begin{theorem}[Main characterization]\label{thm:main-characterization}
Assume $(\X,\rho)$ and $(\Y,\ell)$ are Polish, and $\H$ satisfies the
measurability condition of \citet[Definition~3.3]{bousquet_theory_2021}
with compact $\Theta$ and $h$ continuous in~$\theta$
(cf.\ Lemma~\ref{lem:bridging}).
Then the following are equivalent:
\begin{enumerate}
  \item There exists a (distribution-free) learning rule $\A$ such that for
  every realizable distribution $\bmu$ on $\X\times\Y$,
  the learned predictors $h_n:=\A(S_n)$ satisfy
  \[
    \risk_{\bmu}(h_n)\to 0 \qquad \text{almost surely as } n\to\infty.
  \]
  \item $\H$ admits no infinite non-decreasing-$(\gamma_k)$-Littlestone tree (for
  any non-decreasing gap sequence $(\gamma_k)$ with $\gamma_k\to\infty$).
\end{enumerate}

\end{theorem}

\begin{proof}
$(1) \implies (2)$: Contrapositive: suppose an infinite
  non-decreasing-$\gamma_k$-Littlestone tree exists
  (Definition~\ref{def:monotone-gamma-littlestone}).  By
  Lemma~\ref{lem:bridging}, the tree is automatically realizable
  (Definition~\ref{def:realizable-monotone-gamma-littlestone}).
  Theorem~\ref{thm:lower-bound} then shows that no learner can be
  consistent.
$(2) \implies (1)$: If no such obstruction exists, then the explicit learner constructed in
  Section~\ref{subsec:realizable-learner} is strongly universally Bayes-consistent
